\textbf{Phần 1. Trắc nghiệm lựa chọn}
\setcounterpageref{ex}{0}

\begin{ex} %[3D2V55]
Tính giá trị của phép tính: $7450 - 3280$.
\choice
{\True $4170$}
{$4270$}
{$4130$}
{$4070$}
\loigiai{
$7450 - 3280 = 4170$.
\textbf{Đáp án của câu hỏi là A}
}
\end{ex}

\begin{ex} %[3D2V55]
Một kho hàng có $8500\text{ kg}$ thóc. Người ta đã xuất đi $3650\text{ kg}$. Trong kho còn lại bao nhiêu ki-lô-gam thóc?
\choice
{$4950\text{ kg}$}
{$4800\text{ kg}$}
{\True $4850\text{ kg}$}
{$5150\text{ kg}$}
\loigiai{
$8500 - 3650 = 4850\text{ (kg)}$.
\textbf{Đáp án của câu hỏi là C}
}
\end{ex}

\begin{ex} %[3D2V55]
Tìm số tự nhiên $y$, biết: $9200 - y = 4580$.
\choice
{$4720$}
{\True $4620$}
{$4520$}
{$5620$}
\loigiai{
$y = 9200 - 4580 = 4620$.
\textbf{Đáp án của câu hỏi là B}
}
\end{ex}

\begin{ex} %[3D2V55]
Một cuộn dây cáp dài $5000\text{ m}$. Người ta cắt đi hai lần: lần đầu cắt $1250\text{ m}$, lần sau cắt $1850\text{ m}$. Cuộn dây còn lại dài bao nhiêu mét?
\choice
{$2100\text{ m}$}
{$1800\text{ m}$}
{$2000\text{ m}$}
{\True $1900\text{ m}$}
\loigiai{
Tổng hai lần cắt: $1250 + 1850 = 3100\text{ m}$. Còn lại: $5000 - 3100 = 1900\text{ (m)}$.
\textbf{Đáp án của câu hỏi là D}
}
\end{ex}

\begin{ex} %[3D2V55]
Tìm hiệu của số lớn nhất có 4 chữ số và số tròn trăm lớn nhất có 4 chữ số bé hơn 5000.
\choice
{\True $5099$}
{$4999$}
{$5199$}
{$4899$}
\loigiai{
Số lớn nhất có 4 chữ số là 9999. Số tròn trăm lớn nhất bé hơn 5000 là 4900. Hiệu: $9999 - 4900 = 5099$.
\textbf{Đáp án của câu hỏi là A}
}
\end{ex}

\begin{ex} %[3D2V55]
Một cửa hàng có $6400\ell$ xăng. Sau một tuần bán, cửa hàng còn lại $2750\ell$ xăng. Hỏi cửa hàng đã bán bao nhiêu lít xăng?
\choice
{$3750\ell$}
{$3550\ell$}
{\True $3650\ell$}
{$4650\ell$}
\loigiai{
Số xăng đã bán: $6400 - 2750 = 3650\text{ (\ell)}$.
\textbf{Đáp án của câu hỏi là C}
}
\end{ex}

\begin{ex} %[3D2V55]
Nếu số bị trừ giảm đi 300 đơn vị và số trừ giữ nguyên thì hiệu sẽ
\choice
{Tăng thêm 300 đơn vị}
{\True Giảm đi 300 đơn vị}
{Không thay đổi}
{Giảm đi 600 đơn vị}
\loigiai{
Số bị trừ giảm bao nhiêu đơn vị thì hiệu giảm bấy nhiêu đơn vị.
\textbf{Đáp án của câu hỏi là B}
}
\end{ex}

\begin{ex} %[3D2V55]
Tính nhanh giá trị: $7850 - 2400 - 1850$.
\choice
{$3500$}
{$3700$}
{$3800$}
{\True $3600$}
\loigiai{
Nhóm: $(7850 - 1850) - 2400 = 6000 - 2400 = 3600$.
\textbf{Đáp án của câu hỏi là D}
}
\end{ex}

\begin{ex} %[3D2V55]
Một trang trại thu hoạch được $4500\text{ kg}$ bơ. Đợt một bán $1600\text{ kg}$, đợt hai bán nhiều hơn đợt một $300\text{ kg}$. Trang trại còn lại bao nhiêu ki-lô-gam bơ?
\choice
{\True $1000\text{ kg}$}
{$1200\text{ kg}$}
{$1100\text{ kg}$}
{$900\text{ kg}$}
\loigiai{
Đợt hai bán: $1600 + 300 = 1900\text{ kg}$. Tổng hai đợt bán: $1600 + 1900 = 3500\text{ kg}$. Còn lại: $4500 - 3500 = 1000\text{ (kg)}$.
\textbf{Đáp án của câu hỏi là A}
}
\end{ex}

\begin{ex} %[3D2V55]
Trong phép trừ có hiệu bằng 2450. Nếu tăng số bị trừ thêm 150 đơn vị và tăng số trừ thêm 150 đơn vị thì hiệu mới bằng
\choice
{$2600$}
{$2300$}
{\True $2450$}
{$2750$}
\loigiai{
Khi cả số bị trừ và số trừ cùng tăng một số đơn vị như nhau thì hiệu không thay đổi.
\textbf{Đáp án của câu hỏi là C}
}
\end{ex}

\begin{ex} %[3D2V55]
Quãng đường từ Hà Nội đến Nghệ An dài khoảng $300\text{ km}$. Một bác tài xế đã đi được $185\text{ km}$. Bác tài xế còn phải đi bao nhiêu ki-lô-mét nữa?
\choice
{$125\text{ km}$}
{\True $115\text{ km}$}
{$105\text{ km}$}
{$135\text{ km}$}
\loigiai{
Quãng đường còn lại: $300 - 185 = 115\text{ (km)}$.
\textbf{Đáp án của câu hỏi là B}
}
\end{ex}

\begin{ex} %[3D2V55]
Tìm một số biết rằng lấy số lớn nhất có 4 chữ số trừ đi số đó thì được 3456.
\choice
{$6534$}
{$5543$}
{$6544$}
{\True $6543$}
\loigiai{
Số cần tìm: $9999 - 3456 = 6543$.
\textbf{Đáp án của câu hỏi là D}
}
\end{ex}

\textbf{Phần 2. Trắc nghiệm đúng sai}
\setcounterpageref{ex}{0}

\begin{ex} %[3D2V55]
Xét tính đúng sai của các quy tắc về phép trừ:
\choiceTF
{\True Muốn tìm số trừ ta lấy số bị trừ trừ đi hiệu.}
{Trong phép trừ các số tự nhiên, số trừ luôn lớn hơn số bị trừ.}
{\True Khi cùng thêm vào cả số bị trừ và số trừ một số đơn vị như nhau thì hiệu không thay đổi.}
{$8000 - 2500 = 6500$.}
\loigiai{
\begin{itemize}
    \item Ý a) Đúng: Quy tắc tìm thành phần chưa biết.
    \item Ý b) Sai: Sai vì số bị trừ luôn lớn hơn hoặc bằng số trừ.
    \item Ý c) Đúng: Tính chất của hiệu.
    \item Ý d) Sai: Sai vì $8000 - 2500 = 5500$.
\end{itemize}
\begin{itemize}
    \item a) \textbf{Đúng}.
    \item b) \textbf{Sai}.
    \item c) \textbf{Đúng}.
    \item d) \textbf{Sai}.
\end{itemize}
}
\end{ex}

\begin{ex} %[3D2V55]
Một hồ nước có dung tích $9000\text{ m}^3$ nước, người ta xả bớt nước ra để chống lũ:
\choiceTF
{\True Đợt 1 xả $3200\text{ m}^3$ thì hồ còn lại $5800\text{ m}^3$ nước.}
{\True Đợt 2 xả tiếp $2800\text{ m}^3$ thì sau hai đợt đã xả tổng cộng $6000\text{ m}^3$ nước.}
{Sau hai đợt xả nước trong hồ cạn sạch.}
{\True Lượng nước còn lại trong hồ đúng bằng một phần ba lượng nước ban đầu.}
\loigiai{
\begin{itemize}
    \item Ý a) Đúng: $9000 - 3200 = 5800\text{ m}^3$.
    \item Ý b) Đúng: $3200 + 2800 = 6000\text{ m}^3$.
    \item Ý c) Sai: Sai vì trong hồ còn lại $9000 - 6000 = 3000\text{ m}^3$.
    \item Ý d) Đúng: $3000 : 9000 = \dfrac{1}{3}$.
\end{itemize}
\begin{itemize}
    \item a) \textbf{Đúng}.
    \item b) \textbf{Đúng}.
    \item c) \textbf{Sai}.
    \item d) \textbf{Đúng}.
\end{itemize}
}
\end{ex}

\textbf{Phần 3. Trả lời ngắn}
\setcounterpageref{ex}{0}

\begin{ex} %[3D2V55]
Tính kết quả phép trừ: $8340 - 4565$.
\shortans{3775}
\loigiai{
\textbf{Đáp số:} $3775$\\
$8340 - 4565 = 3775$.
}
\end{ex}

\begin{ex} %[3D2V55]
Tìm số tự nhiên $x$, biết: $7500 - x = 2840$.
\shortans{4660}
\loigiai{
\textbf{Đáp số:} $4660$\\
$x = 7500 - 2840 = 4660$.
}
\end{ex}

\begin{ex} %[3D2V55]
Một trang trại có 5400 con vịt, sau khi bán đi một số con thì còn lại 1850 con. Hỏi trang trại đã bán bao nhiêu con vịt?
\shortans{3550}
\loigiai{
\textbf{Đáp số:} $3550$\\
$5400 - 1850 = 3550$ con.
}
\end{ex}

\begin{ex} %[3D2V55]
Tính nhanh: $6540 - 2300 - 1540$.
\shortans{2700}
\loigiai{
\textbf{Đáp số:} $2700$\\
$(6540 - 1540) - 2300 = 5000 - 2300 = 2700$.
}
\end{ex}

\begin{ex} %[3D2V55]
Một cuộn dây dài $4000\text{ m}$. Lần thứ nhất người ta dùng $1200\text{ m}$, lần thứ hai dùng $1500\text{ m}$. Hỏi cuộn dây còn lại bao nhiêu mét?
\shortans{1300}
\loigiai{
\textbf{Đáp số:} $1300$\\
$4000 - (1200 + 1500) = 1300\text{ (m)}$.
}
\end{ex}

\begin{ex} %[3D2V55]
Hiệu hai số là 1500. Nếu số bị trừ tăng 200 đơn vị và số trừ giảm 100 đơn vị thì hiệu mới bằng bao nhiêu?
\shortans{1800}
\loigiai{
\textbf{Đáp số:} $1800$\\
Hiệu mới: $1500 + 200 + 100 = 1800$.
}
\end{ex}
